% ============================== CI-6 ==============================
\reading{CI-6}{Regulating the Crypto Ecosystem: The Case of Unbacked Crypto Assets}{DEFER}{0}

\begin{cisrcbox}
{\sffamily\bfseries\color{FRMNavy}Why this reading is built in full despite a DEFER tag.}
The strategic reading order gives CI-6 an SRC of 0, meaning no question has been recalled
from it. Current Issues readings rotate annually, so a zero here records the past, not the
next sitting, and the area is not filtered or tiered. Complete coverage is the strategy.
\end{cisrcbox}

\begin{cikeybox}[title={The four objectives, in the spine's own words}]
\begin{enumerate}[label=\textbf{\alph*.}]
\item Define and describe crypto assets, including the categories broadly used by global
financial regulators to classify them.
\item Evaluate the key components within the crypto ecosystem, the potential risks generated
by these components, and potential regulatory responses to address those risks.
\item Identify and describe some of the global approaches to the regulation of unbacked crypto
assets, including the BCBS' proposed treatment of banks' exposures to crypto assets.
\item Examine the considerations and steps introduced by the Bali Fintech Agenda (BFA) for
developing a regulatory framework for crypto assets.
\end{enumerate}
\greyline{Source: ``Regulating the Crypto Ecosystem: The Case of Unbacked Crypto Assets'',
IMF Fintech Notes (September 2022). The class notes label this Reading 105; it is CI-6 on the
2026 spine.}
\end{cikeybox}

% ------------------------------------------------------------------
\lo{CI-6 a}{Define and describe crypto assets, including the categories broadly used by global financial regulators to classify them.}

\begin{cidefbox}[title={Crypto assets}]
\term{Crypto assets} are digital assets that use cryptography and \term{distributed ledger
technology (DLT)} like blockchain to securely record transactions. \term{Blockchain}, the
most common DLT, stores immutable transaction records in linked blocks.\par\vspace{4pt}
Crypto assets can be \term{backed} or \term{unbacked}. Unbacked assets like Bitcoin have no
intrinsic value, rely on supply and demand, and are highly volatile, mainly used for
speculation rather than as true money.
\end{cidefbox}

\subsection*{The categories regulators use}

\begin{center}
\renewcommand{\arraystretch}{1.25}
\small
\begin{tabularx}{\textwidth}{@{}p{2.85cm} >{\raggedright\arraybackslash}X p{3.2cm}@{}}
\toprule
\rowcolor{FRMSoft}
{\sffamily\bfseries\color{FRMNavy}Category} &
{\sffamily\bfseries\color{FRMNavy}Defining characteristics} &
{\sffamily\bfseries\color{FRMNavy}Examples}\\
\midrule
\term{NFT tokens} & Usually centrally issued; right to ownership of a specific product;
collectible and non-substitutable & \\
\addlinespace
\term{Security tokens} & Centrally issued; meets the definition of a security in each
respective jurisdiction, so the definition varies by jurisdiction; within the regulatory
perimeter. Represent rights similar to traditional securities, like profit-sharing or
ownership & \\
\addlinespace
\term{Utility tokens} & Centrally issued; right to a product or service; accepted across
multiple ecosystems; transferable; can be used as a means of exchange. Often within a single
or closed network & Gaming tokens, tokenized store cards\\
\addlinespace
\term{Unbacked crypto assets} & Usually decentralised, sometimes centrally issued; designed
to be used as a means of exchange; limited rights for the token holder; no single issuer to
enforce rights against; transferable. Mainly used for speculation, not payments &
Bitcoin, Ether\\
\addlinespace
\term{Stablecoins} & Designed to be value stable; the stability mechanism can be backing or
collateralisation with a commodity, fiat currency, multiple currencies, crypto assets or
algorithms & Tether, USDC, Binance USD\\
\addlinespace
\term{CBDC} & Centrally issued by a state or central bank; designed to be value stable; the
stability mechanism is usually sovereign fiat currency & \\
\bottomrule
\end{tabularx}
\end{center}
\figcap{Table CI-6.1. The notes' six-column IMF figure, turned on its side so it fits a page
and can be read a row at a time. Two discriminators do most of the work. \emph{Who issues
it}: everything here is centrally issued except unbacked crypto assets, which are usually
decentralised, and that is exactly why there is no issuer to enforce rights against. And
\emph{what stabilises the value}: stablecoins use backing, collateralisation or an algorithm,
while a CBDC uses sovereign fiat, which is the same design goal reached by a different
mechanism. Source: IMF Fintech Notes, Figure 1.}

% ------------------------------------------------------------------
\lo{CI-6 b}{Evaluate the key components within the crypto ecosystem, the potential risks generated by these components, and potential regulatory responses to address those risks.}

\subsection*{1. Issuers}

Decentralised and unbacked crypto assets offer access and anonymity but are hard to regulate.
Disclosures are limited, and networks face risks like cyberattacks. Issuance methods (e.g.,
pre-mining, ICOs) can lead to fraud and manipulation. Marketing materials, including white
papers, often exaggerate benefits and hide risks.

\begin{cikeybox}[title={Regulatory response for issuers}]
\begin{enumerate}[label=\textbf{\alph*)}]
\item \term{White papers}: emphasising clear, accurate, and understandable language with
essential details about the asset and its risks.
\item \term{Notification process}: a mechanism for regulators to be informed about white
paper availability, to monitor developments and risks.
\item \term{Suitability tests}: assessments to determine if crypto assets are appropriate for
retail investors.
\item \term{Marketing}: fair, accurate, and transparent marketing materials aligned with the
white paper.
\item \term{Consumer education}: financial and technology literacy programmes.
\end{enumerate}
\end{cikeybox}

\subsection*{2. Crypto asset exchanges}

Exchanges facilitate buying and selling of crypto assets and may also act as custodians or
market makers by trading from their own inventory.

\begin{center}
\renewcommand{\arraystretch}{1.25}
\begin{tabularx}{\textwidth}{@{}>{\raggedright\arraybackslash}X >{\raggedright\arraybackslash}X@{}}
\toprule
\rowcolor{FRMSoft}
{\sffamily\bfseries\color{FRMNavy}What an exchange does} &
{\sffamily\bfseries\color{FRMRed}Risks it generates}\\
\midrule
\term{Direct retail access}: enables individual traders to engage directly without
intermediaries (no members) &
\term{Custody risks}: cyber, credit, counterparty risk\\
\term{Trade from own inventory}: provides liquidity by allowing exchanges to trade from their
own holdings &
\term{Operational and cyber threats}: hacks leading to bankruptcies and asset loss\\
\term{Custody services}: offers secure storage and faster settlement of crypto assets &
\term{Market manipulation}: high in anonymous, illiquid markets\\
 & \term{PoS risks}: large exchanges can dominate proof-of-stake networks, increasing
centralisation and potential conflicts of interest\\
\bottomrule
\end{tabularx}
\end{center}
\figcap{Table CI-6.2. Read the two columns against each other rather than separately. Each
risk is generated by a function in the left column: custody service creates custody risk,
trading from own inventory puts the exchange on the other side of its customers' trades, and
direct retail access removes the intermediary who would otherwise have been the regulated
party. The concentration of all three functions in one entity is the reading's core concern.}

\term{Challenges in regulation for crypto asset exchanges.} Crypto regulations differ across
countries, with many lacking clear rules. Most focus on lower-risk tokens, allowing exchanges
to avoid oversight by listing unbacked crypto assets. Regulatory gaps remain, especially for
decentralised platforms.

\subsection*{3. Wallets}

Digital wallets store the private keys needed to access crypto assets owned by users. Losing a
wallet means losing the private key and therefore the assets.

\begin{cifigblock}
{\centering
\begin{tikzpicture}[
  font=\sffamily\scriptsize,
  bx/.style={draw=FRMRule, text width=4.55cm, align=left, inner sep=5pt,
             minimum height=1.35cm, rounded corners=1pt},
  ax/.style={font=\sffamily\bfseries\scriptsize\color{FRMNavy}}
]
\node[ax] at (2.6,2.1) {TYPES: who holds the private key};
\node[bx, fill=PaleNavy] (sc) at (2.6,1.15)
  {\textbf{Self-custody}\\[2pt] users control private keys, bearing full responsibility for
   security but risking irreversible loss if keys are lost};
\node[bx, fill=PaleNavy] (cu) at (2.6,-0.55)
  {\textbf{Custodial}\\[2pt] third parties such as exchanges manage keys, reducing user
   burden but introducing risks like cyberattacks or mismanagement};

\node[ax] at (8.6,2.1) {CLASSIFICATION: online or offline};
\node[bx, fill=PaleGold] (hot) at (8.6,1.15)
  {\textbf{Hot wallets}\\[2pt] online, for quick transactions; vulnerable to cyber threats};
\node[bx, fill=PaleGold] (cold) at (8.6,-0.55)
  {\textbf{Cold wallets}\\[2pt] offline, for secure long-term storage; susceptible to
   physical risks such as damage to a USB drive, but preferred by exchanges};

\draw[FRMRule, line width=0.9pt] (5.6,-1.5) -- (5.6,2.3);
\node[font=\sffamily\scriptsize\itshape\color{FRMGrey}, fill=white, inner sep=2pt,
      text width=9.5cm, align=center] at (5.6,-1.95)
  {two independent dimensions: the notes classify wallets on each separately and do not cross
   them into four combined types};
\end{tikzpicture}\par}
\figcap{Figure CI-6.1. The two cuts answer different questions and are easy to swap under
pressure. Self-custody versus custodial is about \emph{who holds the key}, and the risk it
governs is loss versus counterparty failure. Hot versus cold is about \emph{whether the key
is online}, and the risk it governs is cyber versus physical. A custodial wallet can be hot or
cold; the notes do not pair them, so neither does this figure.}
\end{cifigblock}

\begin{cikeybox}[title={Wallet security in crypto regulation: regulatory priorities}]
Regulations aim to protect user assets through segregation, strong cybersecurity, accurate
DLT-based records, AML compliance, prudential norms, cyber insurance, wind-down plans, and
governance to manage conflicts of interest.
\end{cikeybox}

\subsection*{4. Validators, miners, and technology}

\begin{itemize}
\item DLT improves security in financial services through immutable transactions and audit
trails.
\item \term{Private blockchains} offer more control but suffer from fragmentation, limited
interoperability, and centralisation risks.
\item \term{Public blockchains} are resilient to single-point failures but face risks from
user behaviour and governance issues.
\item Most crypto assets lack cross-network interoperability and fungibility across different
blockchains.
\item Regulation aims to be technology-neutral but may clash with innovation goals and broader
policy objectives.
\end{itemize}

\subsection*{5. Regulated financial institutions}

\begin{enumerate}[label=\textbf{\alph*)}]
\item Banks and financial firms face risk from direct exposure to volatile crypto assets and
indirect links (like lending, insurance). Basel proposes risk-based classification, under
which unbacked crypto would attract higher capital requirements.
\item Current regulatory focus is on a conservative approach: capital deductions, stricter
risk weights, and monitoring exposures. Institutions like asset managers, brokers, and
insurers are increasingly offering crypto-related products, increasing systemic risk.
\end{enumerate}

\subsection*{Crypto asset provider functions and market infrastructure}

\begin{enumerate}[label=\textbf{\alph*)}]
\item In crypto markets, single entities handle multiple functions (e.g., custody, clearing),
raising risk. This concentration can lead to disruptions, unlike segmented traditional
markets.
\item Regulators must identify ``systemic'' crypto players (e.g., large stablecoin issuers)
and treat them as \term{Financial Market Infrastructures (FMIs)}.
\end{enumerate}

% ------------------------------------------------------------------
\lo{CI-6 c}{Identify and describe some of the global approaches to the regulation of unbacked crypto assets, including the BCBS' proposed treatment of banks' exposures to crypto assets.}

\subsection*{Global approaches to regulation}

Regulating crypto assets on a national setting may be sufficient in the absence of harmonised
global regulation. Domestic regulation can be bespoke or broad.

\begin{itemize}
\item Examples of different regulatory focus include stablecoins and tokens (EU, UK) and
operational aspects (Japan, Albania). Jurisdictions can be crypto friendly (Switzerland) or
restrictive (Nigeria).
\item Activities relating to crypto assets can bypass domestic rules, especially where
monitoring and enforcement is lacking. As a result, stronger international cooperation among
regulatory authorities is needed to address risks and systemic concerns.
\item Instead of broad bans on crypto assets, regulation should mainly focus on clarifying
rules relating to the economic functions of crypto assets and technical designs.
\end{itemize}

\subsection*{The BCBS proposal}

The BCBS proposes to categorise banks' asset holdings as either (1) lower-risk or (2)
higher-risk (traditional) assets.

\begin{cifigblock}
{\centering
\begin{tikzpicture}[
  font=\sffamily\scriptsize,
  grp/.style={draw=FRMRule, text width=5.4cm, align=left, inner sep=6pt,
              minimum height=1.9cm, rounded corners=1pt}
]
\node[grp, fill=PaleGreen] (lo) at (0,0)
  {\textbf{Group 1: lower-risk assets}\\[3pt]
   risk weights \emph{similar to or lower than} those on traditional assets};
\node[grp, fill=PaleRed, right=8mm of lo] (hi)
  {\textbf{Group 2: higher-risk assets}\\[3pt]
   \textbf{\large 1,250\%} risk weight, applied to \emph{both long and short} positions};
\node[below=3mm of hi.south, font=\sffamily\scriptsize\itshape\color{FRMGrey},
      text width=5.4cm, align=center]
  {a short position does not net the charge away};
\node[below=3mm of lo.south, font=\sffamily\scriptsize\itshape\color{FRMGrey},
      text width=5.4cm, align=center]
  {the crypto wrapper does not by itself raise the weight};
\end{tikzpicture}\par}
\figcap{Figure CI-6.2. Two numbers and one qualifier. A 1,250\% risk weight at an 8\% capital
ratio is a full capital deduction, which is the mechanism behind the ``conservative approach''
described under objective (b). The qualifier is the part most easily dropped: the weight
applies to both long \emph{and} short positions, so a bank cannot offset the charge by
holding the other side.}
\end{cifigblock}

% ------------------------------------------------------------------
\lo{CI-6 d}{Examine the considerations and steps introduced by the Bali Fintech Agenda (BFA) for developing a regulatory framework for crypto assets.}

\begin{cidefbox}[title={Bali Fintech Agenda (BFA)}]
A global regulatory and supervisory framework for regulators that addresses fintech risks,
including crypto assets and technologies.
\end{cidefbox}

It includes the following considerations for crypto asset regulation:

\begin{enumerate}[label=\textbf{\alph*)}]
\item \term{Monitoring} new developments and technologies in the crypto space.
\item \term{Prioritising resources}, with most resources allocated to jurisdictions that have
significant crypto presence, to contain systemic risk.
\item \term{Scope of regulation}, to determine which aspects of crypto activities should be
regulated.
\item \term{Domestic collaboration} and coordination between various regulatory authorities to
address specific risk.
\item \term{Continuous risk assessment} to identify changing risks and provide the strongest
protection to markets and consumers, often through cross-border collaboration.
\end{enumerate}

\begin{cikeybox}[title={Where the five items point}]
Items (a) and (e) are about \emph{watching}, (b) is about \emph{allocating}, (c) is about
\emph{defining the perimeter}, and (d) is about \emph{who talks to whom domestically} while
(e) extends that across borders. The agenda is a framework for regulators to organise
themselves, not a set of rules imposed on issuers.
\end{cikeybox}
